\documentclass[a4paper,twocolumn]{jsarticle}
% a4jでは画像挿入時に問題が起こる場合があるため，a4paperを使用する．

% ===== 基本 =====
\usepackage[
    a4paper,
    left=2.0cm,
    right=2.0cm,
    top=2.0cm,
    bottom=2.0cm
]{geometry}
\usepackage{amsmath,amssymb,mathtools,bm}
\usepackage[dvipdfmx]{graphicx}
\usepackage[dvipdfmx]{xcolor}
\usepackage[dvipdfmx,unicode]{hyperref}
\usepackage{pxjahyper}
\usepackage{titlesec}
\usepackage{wrapfig}
\usepackage{multicol}
\usepackage{enumitem}
\usepackage{listings}
\usepackage{float}
\usepackage{caption}
\usepackage{subcaption}

% ===== TikZ・図 =====
\usepackage{tikz}
\usetikzlibrary{positioning}
\usepackage{pgfplots}
\pgfplotsset{compat=1.18}

% ===== 化学 =====
\usepackage[version=4]{mhchem}
\usepackage{chemfig}

% ===== 回路 =====
\usepackage{circuitikz}

% ===== アルゴリズム =====
\usepackage{algorithm}
\usepackage{algpseudocode}

% ===== ダミーテキスト =====
% 発展例で図の回り込みを確認するために使用する．実際の文書では不要．
\usepackage{lipsum}

% ===== 数学補助 =====
% \numberwithin{equation}{section}
\renewcommand{\theequation}{\arabic{equation}}

\DeclareMathOperator{\sgn}{sgn}
\DeclareMathOperator{\tr}{tr}
\DeclareMathOperator{\diag}{diag}

\newcommand{\dif}{\mathrm{d}}
\newcommand{\ee}{\mathrm{e}}
\newcommand{\ii}{\mathrm{i}}

\newcommand{\pd}[2]{\frac{\partial #1}{\partial #2}}
\newcommand{\pdd}[2]{\frac{\partial^2 #1}{\partial #2^2}}
\newcommand{\dd}[2]{\frac{\mathrm{d} #1}{\mathrm{d} #2}}

\newcommand{\vect}[1]{\bm{#1}}
\newcommand{\tensor}[1]{\bm{#1}}

% 画像がまだ配置されていない場合も，文書全体はコンパイルできるようにする．
% 指定した画像が存在すれば，通常どおりその画像を表示する．
\newcommand{\practiceimage}[2][\linewidth]{%
  \IfFileExists{#2}{%
    \includegraphics[width=#1]{#2}%
  }{%
    \fbox{%
      \parbox[c][28mm][c]{\dimexpr#1-2\fboxsep-2\fboxrule\relax}{%
        \centering\small
        画像ファイル\\未配置%
      }%
    }%
  }%
}

% ===== 見出し設定 =====
\titleformat{\section}[block]
  {\normalfont\Large\bfseries}
  {\thesection}{1em}{}

\titleformat{\subsection}[block]
  {\normalfont\large\bfseries}
  {\thesubsection}{1em}{}

\titleformat{\subsubsection}[block]
  {\normalfont\normalsize\bfseries}
  {\thesubsubsection}{1em}{}

\titlespacing*{\section}{0pt}{2.0ex plus 0.5ex minus 0.2ex}{0.8ex}
\titlespacing*{\subsection}{0pt}{1.5ex plus 0.4ex minus 0.2ex}{0.5ex}

\lstset{
    language=C,
    basicstyle=\ttfamily\small,
    numbers=left,
    numberstyle=\tiny\color{gray},
    frame=single,
    rulecolor=\color{black},
    backgroundcolor=\color{gray!10},
    breaklines=true,
    showstringspaces=false,
    columns=fullflexible,
    keywordstyle=\color{blue},
    commentstyle=\color{green!50!black},
    stringstyle=\color{red},
    tabsize=4
}

\setlist[itemize]{topsep=4pt,itemsep=2pt,leftmargin=*}
\setlist[enumerate]{topsep=4pt,itemsep=2pt,leftmargin=*}
\captionsetup{font=small}
\setlength{\columnsep}{8mm}
\setlength{\emergencystretch}{1em}
\raggedbottom

\title{人物識別システムの設計と実装}
\author{東海大学 情報理工学部\\
コンピュータ応用工学科\\
学籍番号：XXXXXXXX\\
徐 海杰}
\date{\today}

\begin{document}

\maketitle

\begin{abstract}
本研究では，YOLOv5，DeepSORT，ReIDおよび人物属性認識を用いた人物識別システムを提案する．
監視カメラ映像から人物を検出・追跡し，同一人物を識別することで入退室管理を実現した．
実験ではシステムの認識性能およびリアルタイム性能を評価した．
\end{abstract}

\section{はじめに}
近年，監視カメラを利用した人物認識技術は，
防犯や入退室管理など様々な分野で利用されている．
本研究では，YOLOv5，DeepSORT，ReIDおよび
人物属性認識を組み合わせた人物識別システムを提案する．

\section{システム構成}
本システムは人物検出，人物追跡，人物識別，
データベース管理の4つのモジュールから構成される．

\section{人物検出}
人物検出にはYOLOv5を用いた．
YOLOv5はリアルタイム物体検出が可能であり，
人物クラスのみを検出対象とした．

\begin{verbatim}
\begin{abstract}
ここにAbstractを記述します．
\end{abstract}
\end{verbatim}

実際に，このコードを\verb|\maketitle|の後へ挿入してみましょう．
この解答例では，タイトルの直後にダミーテキストのAbstractを挿入しています．


\section{人物追跡}
人物追跡にはDeepSORTを用いた．
DeepSORTは検出結果と特徴量を利用して，
各人物にTrack IDを割り当てる．

\subsection*{インライン数式}
円の面積は $S=\pi r^2$ で表される．

オイラーの公式は $e^{i\pi}+1=0$ である．

関数 $f(x)=x^2+2x+1$ を考える．

\subsection*{IoU（YOLO）}
\begin{equation}
IoU=\frac{|A\cap B|}{|A\cup B|}
\end{equation}

\subsection*{ReID}
\begin{equation}
d(\mathbf{x},\mathbf{y})
=
\sqrt{\sum_{i=1}^{n}(x_i-y_i)^2}
\end{equation}

\subsection*{Cosine Similarity}
\begin{equation}
\cos\theta=
\frac{\mathbf{x}\cdot\mathbf{y}}
{\|\mathbf{x}\|\|\mathbf{y}\|}
\end{equation}

ディスプレイ数式は，\verb|\[ ... \]|で囲むと番号なしで表示できます．
番号を付けたい場合は，\verb|equation|環境を使用します．
実際に，インライン数式とディスプレイ数式をそれぞれ使って，
自分で数式を挿入してみましょう．


\section{人物識別}
人物識別にはReID特徴量と
Paddle Attributeによる属性認識を組み合わせた．


\begin{figure}[H]
    \centering
    \practiceimage[3cm]{figures/LaTeX_practice_figure_1.jpg}
    \caption{LaTeXで挿入した図の例}
    \label{fig:practice}
\end{figure}

画像の挿入は，画像ファイルの場所や拡張子，パッケージ設定などが原因で
エラーになることがあります．
エラーが出ても慌てず，エラーメッセージを読みながら一つずつ原因を確認しましょう．
では，自分で画像を挿入してみましょう．

\subsection*{解答例}
\begin{figure}[H]
    \centering
    \practiceimage[3cm]{figures/LaTeX_practice_figure_2.jpg}
    \caption{自分で挿入した図の例}
    \label{fig:tokai}
\end{figure}

\section{実験}
実験環境および認識精度について評価した

\begin{table}[H]
\centering
\caption{人物識別精度}
\label{tab:accuracy}

\begin{tabular}{|c|c|}
\hline
手法 & Accuracy(\%)\\
\hline
YOLOv5 & 96.2\\
DeepSORT & 94.8\\
提案手法 & 97.5\\
\hline
\end{tabular}

\end{table}

\begin{table}[H]
    \centering
    \caption{各手法の認識性能}
    \label{tab:result}

    \resizebox{\columnwidth}{!}{%
    \begin{tabular}{|c|c|c|c|c|}
        \hline
        手法 & Precision(\%) & Recall(\%) & Accuracy(\%) & FPS \\
        \hline
        YOLOv5 & 95.6 & 94.8 & 95.2 & 32 \\
        \hline
        YOLOv5+DeepSORT & 96.4 & 95.9 & 96.1 & 28 \\
        \hline
        提案手法 & 97.5 & 97.0 & 97.2 & 24 \\
        \hline
    \end{tabular}%
    }
\end{table}

\section{まとめ}
本研究ではYOLOv5，DeepSORT，
ReIDおよび属性認識を組み合わせた人物識別システムを構築した．

\begin{itemize}
    \item LaTeX
    \item Python
    \item GitHub
\end{itemize}

では，自分で箇条書きを作成してみましょう．

\subsection*{解答例}
\begin{itemize}
    \item 数学
    \item 物理
    \item 化学
    \item 生物
    \item コンピュータサイエンス
\end{itemize}

\begin{lstlisting}
#include <stdio.h>

int main(void)
{
    for (int i = 0; i < 10; i++) {
        printf("%d\n", i);
    }

    return 0;
}
\end{lstlisting}

では，自分でソースコードを挿入してみましょう．

\subsection*{解答例}
\begin{lstlisting}
#include <stdio.h>

typedef struct {
    char name[20];
    int age;
} Student;

int main(void)
{
    Student s = {"Alice", 20};

    printf("Name : %s\n", s.name);
    printf("Age  : %d\n", s.age);

    return 0;
}
\end{lstlisting}

\section*{演習9：URLの挿入}
URLを挿入することもできます．
OpenAIのホームページはこちらです．

\url{https://openai.com}

では，自分で任意のURLを挿入してみましょう．

\subsection*{解答例}
慎先生のresearchmapはこちらです．

\url{https://researchmap.jp/sgshin}

\section*{演習10：参考文献}
卒業論文や学術論文では，他者の研究や論文を引用した場合，
必ず参考文献を明記する必要があります．

LaTeXでは，一般的にBibTeXを用いて参考文献を管理します．
参考文献を追加する手順は，大きく分けて次の3つです．

\subsection*{手順1：\texttt{.bib}ファイルを作成する}
まず，\texttt{references.bib}という名前のファイルを作成します．
このファイルには，引用する論文や書籍の情報を記述します．
例えば，次のように記述します．

\begin{lstlisting}[language={}]
@inproceedings{he2016resnet,
  author    = {Kaiming He and others},
  title     = {Deep Residual Learning for Image Recognition},
  booktitle = {CVPR},
  year      = {2016}
}
\end{lstlisting}

Google ScholarやIEEE Xploreでは，「BibTeX」を選択するだけで
この形式を取得できるため，自分で一から入力する必要はほとんどありません．

\subsection*{手順2：本文中で引用する}
引用したい箇所に，次のように記述します．

\begin{verbatim}
\cite{he2016resnet}
\end{verbatim}

例えば，

\begin{verbatim}
深層学習ではResNetが広く利用されている
\cite{he2016resnet}．
\end{verbatim}

と書くと，「深層学習ではResNetが広く利用されている[1]．」のように，
引用番号が自動で付与されます．

\subsection*{手順3：参考文献一覧を表示する}
文書の末尾へ，次の2行を追加します．

\begin{verbatim}
\bibliographystyle{ieeetr}
\bibliography{references}
\end{verbatim}

ここで，

\begin{itemize}
    \item \texttt{ieeetr}：参考文献の表示形式
    \item \texttt{references}：作成した\texttt{references.bib}の指定
\end{itemize}

を表しています．
実際に手順1から手順3までを行うと，本文では次のように表示されます．

「深層学習ではResNetが広く利用されている\cite{he2016resnet}．」

卒業論文では，参考文献の管理にBibTeXを用いることが一般的です．
今回は演習を省略しますが，研究を進める際には必ず利用する機能です．

\subsection*{発展：2段組で大きな図を挿入する}
卒業論文では，本文を2段組にし，図やグラフのみページ全体の幅で
表示することがあります．

その場合は，\verb|figure|の代わりに\verb|figure*|環境を使用します．
図がページ上部へ自動配置されることを確認するため，後ろにダミーテキストを入れています．

\begin{lstlisting}[language={[LaTeX]TeX}]
\begin{figure*}[t]
    \centering
    \includegraphics[width=0.8\textwidth]{graph.pdf}
    \caption{実験結果}
    \label{fig:result}
\end{figure*}
\end{lstlisting}

\begin{figure*}[t]
    \centering
    \begin{subfigure}{0.47\textwidth}
        \centering
        \practiceimage[\linewidth]{figures/LaTeX_practice_figure_1.jpg}
        \caption{画像1}
        \label{fig:image1}
    \end{subfigure}
    \hfill
    \begin{subfigure}{0.47\textwidth}
        \centering
        \practiceimage[\linewidth]{figures/LaTeX_practice_figure_2.jpg}
        \caption{画像2}
        \label{fig:image2}
    \end{subfigure}
    \caption{画像を横並びにした例}
    \label{fig:images}
\end{figure*}

\lipsum[2-7]

% 参考文献一覧は，文書の最後に一度だけ出力する．
\clearpage
\bibliographystyle{ieeetr}
\bibliography{references}

\end{document}
